\documentclass[10pt]{article}
\usepackage[margin=0.9in]{geometry}
\usepackage{amsmath,amssymb,amsthm,booktabs}
\usepackage[hidelinks]{hyperref}
\newtheorem{lemma}{Lemma}
\newtheorem{theorem}{Theorem}
\title{Unbounded speed replacements for seventeen lonely runners}
\author{Claude Make No Mistakes}
\date{26 November 2026}
\begin{document}
\maketitle

\begin{abstract}
Replace up to six of the speeds \(1,\ldots,16\) by distinct
integers greater than 16. We show by exact interval computation
that every nontrivial replacement set has loneliness strictly
greater than \(1/17\). The inserted speeds have no imposed upper
bound. A density bound on each surviving time interval determines
a finite range for the smallest remaining inserted speed.
\end{abstract}

\section{An interval bound}
Let \(S\) be a nonempty finite set of positive integer speeds and put
\[
 F_S(t)=\min_{v\in S}\|vt\|,\qquad
 U(S)=\{t\in[0,1/2]:F_S(t)\geq 1/17\}.
\]
Here \(\|x\|\) denotes distance to the nearest integer.
Periodicity and reflection allow us to restrict to \([0,1/2]\).
The region \(U(S)\) is a finite union of closed intervals with
rational endpoints, possibly singletons. If it has positive
length, then some time satisfies \(F_S(t)>1/17\), since equality
for any one speed occurs at finitely many times.

Zhang bounds the measure that an inserted runner can cover in an
uncovered interval in the study of two-entry modifications
\cite[\S6.1]{zhang}. We use the sharp single-interval discrepancy
and apply it recursively, updating the interval after every choice.

\begin{lemma}
Suppose \(I\subseteq U(S)\) is an interval of length \(L>0\).
Let \(1\leq r\leq8\), and let \(w_1<\cdots<w_r\) be positive
integer speeds such that \(\max_t F_{S\cup\{w_1,\ldots,w_r\}}(t)
\leq1/17\). If \(r\geq2\), then
\[
 w_1\ \leq\ \frac{30r}{17(17-2r)L}.
\]
If \(r=1\), then \(w_1\leq2/(17L)\).
\end{lemma}

\begin{proof}
Set \(\alpha=2/17\). For a single speed \(w\), the closed set of
times with \(\|wt\|\leq1/17\) occupies an \(\alpha\)-fraction
of each period of length \(1/w\). Write \(wL=m+\theta\), with
\(m\) integral and \(0\leq\theta<1\). In the \(wt\) coordinate,
the \(m\) full periods contribute \(m\alpha\). The remaining
interval projects injectively to the period circle except at
endpoints, so its contribution is at most \(\min(\theta,\alpha)\).
Since \(\min(\theta,\alpha)-\alpha\theta\leq\alpha(1-\alpha)\),
the bad portion of \(I\) has length at most
\(\alpha L+\alpha(1-\alpha)/w\).

At every time in \(I\) other than the finitely many retained-speed
equality times, at least one of the new speeds must be bad.
The union bound gives
\[
 L\ \leq\ r\alpha L+\alpha(1-\alpha)\sum_{i=1}^r\frac1{w_i}
 \ \leq\ r\alpha L+\frac{r\alpha(1-\alpha)}{w_1}.
\]
Solving yields the claimed bound. If \(r=1\), the closed bad
set covers a dense subset of \(I\), hence all of \(I\).
Connectedness forces \(I\) into one bad interval, of length
\(2/(17w_1)\).
\end{proof}

\section{Complete search with no prescribed speed cap}
\begin{theorem}
Let \(V\) be a set of 16 distinct positive integers. If at most
six elements of \(V\) exceed 16, then
\(\max_t F_V(t)\geq1/17\), with equality precisely for
\(V=\{1,\ldots,16\}\).
\end{theorem}

\begin{proof}
For the baseline, the seventeen points \(0,t,\ldots,16t\)
on the unit circle have a gap of length at most \(1/17\).
Subtracting its endpoints gives a speed at distance at most
\(1/17\) from the stationary runner. At \(t=1/17\), equality
holds.

For each \(r\in\{1,\ldots,6\}\), choose the \(r\) removed speeds
from \(\{1,\ldots,16\}\) and start from the retained speeds.
We search for all completions with loneliness at most \(1/17\).
If there are \(k\) inserted speeds still to choose, use a
longest interval in the current \(U(S)\) and the lemma to bound
the least remaining speed. Branch over that finite range,
respecting the increasing order and the lower bound 17.
Intersect \(U(S)\) with the chosen speed's good intervals,
whose endpoints at time \(t\) are
\[
 \frac{17j+1}{17v}\qquad\text{and}\qquad
 \frac{17j+16}{17v}\qquad(j\in\mathbb Z).
\]
At the last choice we first require every positive interval
to fit in a single closed bad arc of that speed. A full tuple
with empty \(U\) is a counterexample. A nonempty, finite \(U\)
gives a tight instance.

At any proper prefix with no positive interval, only finitely
many rational times remain. An additional speed divisible by
all of their denominators rules out all those times, giving
a strict counterexample after padding with distinct larger
speeds. This check is included before the interval bound is
applied. The completed search encountered no such prefix.
All retained sets were examined and no full completion
with loneliness at most \(1/17\) remained.
\end{proof}

\section{Certificates and scope}
The archives include one record per retained set and store the
explored counts and maximum computed cutoffs at each depth.
There are 14,892 choices across one to six replacements.
Both search implementations keep interval endpoints exact.
The native implementation stops on any bound exceeding its
arithmetic range. No such bound was reached in the recorded
search. The separate verifier enumerates the expected retained
sets and recomputes selected rows using rational arithmetic
and a direct last-speed intersection. It also checks the
tight baseline and a strict-cover control.

\begin{center}
\begin{tabular}{rrrr}
\toprule
Replacements & Retained sets & Speed trials & Largest cutoff \\
\midrule
1 & 16 & 72 & 32 \\
2 & 120 & 1,899 & 64 \\
3 & 560 & 44,575 & 98 \\
4 & 1,820 & 826,808 & 141 \\
5 & 4,368 & 14,879,151 & 175 \\
6 & 8,008 & 346,744,936 & 306 \\
\bottomrule
\end{tabular}
\end{center}

Goddyn and Wong study simultaneous multiplicative speedups
\cite{goddyn}. Zhang classifies arbitrary one-entry replacements
and reports a bounded two-entry census \cite{zhang}. The
present range keeps the runner count fixed and allows up to
six arbitrary replacements above 16. Any counterexample to
the unrestricted seventeen-runner conjecture would therefore
require at least seven speeds above 16 in a primitive integer
speed set. The reduction to positive integer speeds is standard
\cite{tao}, and division by their common divisor preserves
the maximum loneliness.

\begin{thebibliography}{9}
\bibitem{goddyn} L. Goddyn and E. B. Wong.
Tight instances of the lonely runner.
\emph{Integers} 6 (2006), A38.
\bibitem{tao} T. Tao. Some remarks on the lonely runner conjecture.
\emph{Contributions to Discrete Mathematics} 13 (2018), 1--31.
\href{https://arxiv.org/abs/1701.02048}{arXiv:1701.02048}.
\bibitem{zhang} Y. Zhang. Tight instances of the Lonely Runner
Conjecture: one-entry modifications and the growth of speeds.
\href{https://arxiv.org/html/2608.13599v3}{arXiv:2608.13599v3}, 2026.
\end{thebibliography}
\end{document}
